\chapter{Sprint K --- Backend Bù hàng/Return: redesign + wizard SO 0đ + hook giao thực tế (2026-07-17)}

\section{Bối cảnh}

BA giao chức năng backend \textbf{Bù hàng / Đổi trả} (Google Sheet ``Internal ERP Master Plan\_Update'',
tab ``1. Model/Field''). Tham chiếu v14 \texttt{co\_franchise\_return} (allocation + FIFO + floor UoM +
SO 0đ wizard). Sprint chia 2 phần, đóng sổ CHUNG 1 lượt (1 chapter, 1 commit) theo quyết định user:
\begin{itemize}
  \item \textbf{K1 --- Foundation}: redesign \texttt{wujia\_portal\_return} từ multi-line sang
    single-product + khối phê duyệt + theo dõi phân bổ (compute) + workflow state + backend admin
    + portal single-product + migration \texttt{19.0.2.0.0}.
  \item \textbf{K2 --- Xử lý bù hàng}: wizard ``Xử lý bù hàng'' tạo SO 0đ + phân bổ FIFO, và hook
    \texttt{stock.picking} cập nhật số đã-bù từ giao thực tế.
\end{itemize}

\textbf{8 điểm BA đã chốt (2026-07-17):} (1) state \texttt{draft$\to$submitted$\to$reviewing$\to$approved$\to$processing$\to$done}
+ \texttt{rejected}/\texttt{cancelled}; (2) giá SO bù \texttt{price\_unit=0}; (3) ĐVT giao quy đổi qua
engine UoM Odoo (\texttt{\_compute\_quantity}) cùng nhóm --- accumulate dùng hệ số \texttt{unit\_qty} tay
(bịch $\neq$ kg khác nhóm); (4a) ``đã bù'' \textbf{chỉ} tính từ giao thực tế (\texttt{stock.picking} done),
cộng dồn từng đợt; (4b) kho giao thiếu không backorder $\to$ tự hoàn phần chênh, bù tiếp lần sau;
(5) \textbf{1 SO / 1 cửa hàng / 1 lần xử lý}; (6) FIFO tự đề xuất SL, HQ chỉnh tay; (7) từ chối sau khi
bù một phần $\to$ SO giữ nguyên; (8) \texttt{compensation\_unit\_qty} có default trên product.

\section{K1 --- Foundation}

\begin{itemize}
  \item \textbf{\texttt{product.product}} (extend): \texttt{compensation\_enabled}, \texttt{compensation\_policy}
    (exact/accumulate), \texttt{compensation\_claim\_uom\_id}, \texttt{compensation\_product\_id},
    \texttt{compensation\_delivery\_uom\_id}, \texttt{compensation\_unit\_qty} (default gợi ý khi duyệt);
    constraint conditional theo policy. Tab ``Bù hàng'' ở product form (\texttt{group\_return\_manager}).
  \item \textbf{\texttt{wujia.return.issue.type}} (rename từ \texttt{error.type}): giữ ormcache.
  \item \textbf{\texttt{wujia.return.request}} redesign + \texttt{mail.thread}: header single-product
    (\texttt{sale\_order\_line\_id}$\to$\texttt{product\_id}/\texttt{request\_qty}, thay \texttt{line\_ids});
    khối duyệt (\texttt{approved\_qty}/\texttt{approved\_uom\_id} + snapshot bù); theo dõi phân bổ
    (\texttt{allocated\_qty}/\texttt{compensated\_qty}/\texttt{unallocated\_qty}/\texttt{remaining\_qty}
    compute qua \texttt{\_read\_group} --- 1 query, perf 1500 user); workflow
    \texttt{action\_submit/start\_review/approve/reject/cancel/mark\_done}.
  \item \textbf{\texttt{wujia.compensation.allocation}}: bản ghi phân bổ (1 request $\times$ 1 dòng SO bù),
    \texttt{allocated\_qty}/\texttt{delivered\_qty}/\texttt{released\_qty}/\texttt{open\_qty} theo
    \texttt{allocation\_uom\_id} (= \texttt{approved\_uom\_id}); constraint SQL+Python
    (\texttt{delivered+released $\le$ allocated}; $\sum$(allocated$-$released) $\le$ approved\_qty).
  \item \textbf{\texttt{sale.order.is\_return\_order}} (thêm ở \texttt{wujia\_sale}).
  \item \textbf{Portal single-product}: controller \texttt{\_parse\_payload} + rewrite form/detail/list
    (desktop về single-product như mobile). \textbf{Migration \texttt{19.0.2.0.0}} pre/post: drop bảng
    line cũ, rename cột, rename \texttt{ir.model} error$\to$issue, dọn ACL/rule.
\end{itemize}

\section{K2 --- Xử lý bù hàng (SO 0đ) + hook giao thực tế}

\textbf{Wizard 3 TransientModel} \texttt{wujia.compensation.process.wizard(.group)(.line)} +
\texttt{ir.actions.act\_window} bind vào list (menu Actions, \texttt{group\_ids} = manager). 8 bước:
\begin{enumerate}
  \item Validate: chỉ request \texttt{approved}/\texttt{processing} + \texttt{resolution\_type='compensation'}
    + đủ snapshot + \texttt{unallocated\_qty > 0}.
  \item Group by \texttt{(franchise, comp\_product, approved\_uom, delivery\_uom, policy, unit\_qty)} ---
    nhiều group cùng franchise $\to$ CÙNG 1 SO (điểm 5).
  \item FIFO \texttt{(approved\_date, request\_date, id)}.
  \item \texttt{total\_claim} = $\sum$ \texttt{unallocated\_qty} quy về ĐVT quyền lợi.
  \item Đề xuất SL giao: \texttt{exact} = convert (round DOWN, không bù vượt); \texttt{accumulate} =
    \texttt{floor(total/unit\_qty)}, phần lẻ để lại.
  \item HQ chỉnh \texttt{delivery\_qty} (điểm 6).
  \item Confirm re-validate: \texttt{SELECT ... FOR UPDATE NOWAIT} + đọc lại \texttt{unallocated\_qty}
    \textbf{và state} $\to$ lệch snapshot / bị huỷ $\to$ raise ``Dữ liệu đã thay đổi''.
  \item Tạo 1 SO/franchise (\texttt{is\_return\_order=True}, \texttt{price\_unit=0}, draft) + từng dòng SO
    (gắn allocation ngay, không phụ thuộc thứ tự) + set request \texttt{processing}.
\end{enumerate}

\textbf{Hook \texttt{stock.picking.\_action\_done}} (chỉ phiếu \texttt{outgoing}): với move thuộc SO
\texttt{is\_return\_order}, cộng dồn SL thực giao vào \texttt{allocation.delivered\_qty} (điểm 4a ---
backorder tự cộng dồn); giao xong không backorder mà còn open $\to$ \texttt{released\_qty} hoàn phần chênh
(điểm 4b); request \texttt{processing}$\to$\texttt{done} khi bù đủ. \texttt{compensated\_qty} \textbf{chỉ}
tăng ở đây, không từ lúc tạo/confirm SO.

\section{Hardening (adversarial review 5-lens)}

Review đối kháng (5 lens $\to$ skeptic verify từng finding, 19 confirmed $\to$ 8 root) tìm ra và đã sửa:
\begin{itemize}
  \item \textbf{(blocker) Group key thiếu \texttt{approved\_uom\_id}} $\to$ 2 request khác ĐVT quyền lợi
    (kg/g) gộp 1 dòng SO, hook \texttt{\_fill\_allocations} trộn sai ĐVT. Fix: thêm
    \texttt{approved\_uom\_id} vào key $\to$ mỗi group đồng nhất ĐVT.
  \item \textbf{(blocker) Exact double-UP rounding} $\to$ \texttt{covered > total\_claim} raise ngay trên
    SL wizard tự đề xuất. Fix: quy đổi round \textbf{DOWN} 2 chiều (\texttt{\_claim\_to\_delivery} +
    \texttt{\_delivery\_to\_claim}), covered lấy từ back-conversion.
  \item \textbf{(high) Portal ACL} \texttt{wujia.return.request} cho \texttt{group\_portal}
    \texttt{write/create=1} $\to$ user portal có thể tự duyệt (set \texttt{state=approved},
    \texttt{approved\_qty}) qua RPC. Fix: hạ ACL về \textbf{read-only} (controller portal đã dùng
    \texttt{sudo()} cho mọi create/write $\to$ không ảnh hưởng luồng portal).
  \item \textbf{(high) Re-validate không kiểm state} $\to$ request bị huỷ giữa chừng vẫn tạo allocation
    (\texttt{unallocated\_qty} không đổi khi chỉ đổi state). Fix: kiểm \texttt{state} trong
    \texttt{\_lock\_and\_revalidate}.
  \item \textbf{(medium) Odoo 19 \texttt{\_compute\_quantity} bỏ kiểm nhóm} (\texttt{category\_id} $\to$ cây
    \texttt{relative\_uom\_id}). Fix: tự guard bằng gốc cây \texttt{\_uom\_root}, khác nhóm $\to$ raise.
  \item \textbf{(medium/perf) Index + hook outgoing-only}: thêm \texttt{index=True} cho
    \texttt{allocation.sale\_order\_line\_id} (hot-path search); hook chỉ chạy phiếu \texttt{outgoing}
    (early-exit phiếu nhập/nội bộ 1500 user + không đếm nhầm move phiếu trả hàng).
\end{itemize}

\section{Gì đã làm (nghiệp vụ)}

Hoàn thiện toàn bộ quy trình \textbf{bù hàng} phía văn phòng (HQ): cửa hàng gửi yêu cầu bù (1 sản phẩm/phiếu,
kèm ảnh minh chứng) qua portal; HQ tiếp nhận, phê duyệt số lượng và chọn sản phẩm/đơn vị bù. Khi HQ bấm
``Xử lý bù hàng'' cho nhiều yêu cầu, hệ thống tự gom theo cửa hàng thành \textbf{một đơn bù 0đ}, đề xuất số
lượng giao theo FIFO (ưu tiên yêu cầu cũ) --- HQ chỉnh tay được. Số ``đã bù'' chỉ được ghi nhận khi hàng
\textbf{thực sự giao tới cửa hàng}; giao làm nhiều đợt thì cộng dồn, giao thiếu thì phần còn lại tự động
chờ bù đợt sau. Yêu cầu tự chuyển ``Đã xử lý'' khi bù đủ.

\section{Lý do/Trade-off}

Tách ĐVT ở 2 phía là quyết định cốt lõi: \texttt{allocation} lưu theo \textbf{ĐVT quyền lợi} (để so trực
tiếp với \texttt{approved\_qty}), còn dòng SO lưu \textbf{ĐVT giao} --- quy đổi qua engine UoM Odoo (exact)
hoặc hệ số \texttt{unit\_qty} tay (accumulate, vì ``bịch $\to$ kg'' khác nhóm đơn vị nên Odoo không tự đổi).
SO bù để ở \texttt{draft} cho HQ review rồi tự confirm (không auto-confirm). Rounding chọn \textbf{DOWN} để
không bao giờ bù vượt quyền lợi. Bảo mật: thay vì viết \texttt{write()} override phức tạp, tận dụng việc
controller portal đã \texttt{sudo()} toàn bộ $\to$ chỉ cần hạ ACL portal về read-only là đóng lỗ tự-duyệt.

\section{Verify}

\texttt{-u wujia\_portal\_return,wujia\_sale}: RC=0, 0 ERROR log, migration pre/post sạch. ORM test
\textbf{46/46} (13 kịch bản: exact/accumulate, HQ chỉnh, giao đủ/backorder/thiếu, concurrency, over-approved,
FIFO 2 request, 1 SO đa dòng, cross-UoM tách group, huỷ giữa chừng, exact lẻ không over-cover, price=0
sau confirm). \texttt{test\_sprint9} 7/7. Portal + backend routes HTTP 200. Manifest:
\texttt{wujia\_portal\_return 19.0.2.0.0} (K1+K2 cùng release) / \texttt{wujia\_sale 19.0.4.1.0} /
\texttt{wujia\_portal\_base 19.0.5.14.0}. Deploy prod: \texttt{-u wujia\_portal\_return wujia\_sale}
(module đã cài $\to$ auto-upgrade chạy migration; KHÔNG cần \texttt{-i}).

\section{Portal wire --- Tiến độ bù cho cửa hàng (2026-07-17)}

Phần bổ sung đóng nốt Sprint K: toàn bộ K1+K2 chạy ở backend HQ nên cửa hàng vào
\texttt{/portal/return/<id>} chỉ thấy thông tin mình gửi + lý do từ chối, không thấy kết quả xử lý bù.
Sprint này \textbf{kéo dữ liệu bù đã có sẵn} (compute từ allocation) ra hiển thị cho cửa hàng ---
\textbf{read-only, không đổi schema}.

\begin{itemize}
  \item \textbf{Module}: \texttt{wujia\_portal\_return 19.0.2.0.0 $\to$ 19.0.2.1.0} (UI-only).
  \item \textbf{Controller} (\texttt{controllers/portal.py}): thêm map \texttt{RESOLUTION\_LABELS} +
    \texttt{COMPENSATION\_STATUS\_LABELS}; helper \texttt{\_build\_compensation\_ctx(rr)} (trả \texttt{None}
    khi HQ chưa chốt phương án $\to$ ẩn card, không lộ số 0) + \texttt{\_so\_delivery\_label(so)} (đếm
    phiếu giao done/tổng, guard \texttt{'picking\_ids' in so.\_fields}). Truyền \texttt{comp}/\texttt{resolution\_labels}
    vào detail, \texttt{comp\_status\_labels} vào list.
  \item \textbf{Template detail} (desktop + mobile): card ``Phương án xử lý \& Tiến độ bù'' ---
    phương án, sản phẩm bù, SL duyệt bù, badge tình trạng, progress bar (\texttt{--wujia-primary}),
    ``Đã bù X / Còn thiếu Y'', danh sách đơn bù + trạng thái giao; bổ sung \texttt{approval\_note}
    (phản hồi duyệt HQ) cạnh \texttt{reject\_reason}.
  \item \textbf{Template list} + CSS: badge \texttt{compensation\_status} cho dòng bù; progress-bar
    style (\texttt{.wujia-return-progress}) + block mobile trong \texttt{portal\_return.css}.
\end{itemize}

\textbf{Quyết định:} không làm nút ``Đặt hàng bù'' cho cửa hàng --- kiến trúc bù hàng do HQ chạy wizard
tạo SO 0đ, cửa hàng không tự đặt; cái cửa hàng cần là \emph{thấy tiến độ}. Nút ``xác nhận đã nhận bù''
= add-on sprint sau nếu BA yêu cầu.

\textbf{Verify:} \texttt{-u wujia\_portal\_return} RC=0, không ParseError. Render có xác thực
\texttt{/portal/return/<id>} HTTP 200 hiển thị đúng progress 30\%, badge ``Đang bù'', danh sách đơn bù,
phản hồi HQ; list + form 200; \texttt{test\_sprint9} 7/7; smoke 8 trang portal 200. Deploy: UI-only,
auto-deploy (bump manifest), \textbf{KHÔNG} cần \texttt{-i}, không schema.
